\documentclass{osucourse}
\course{1135}

%% Not in the university's course data, so these come
%% from the published sheet this file was imported from.
\SetCourseField{title}{Number and Operations for Teachers}
\SetCourseField{credits}{5 credits}
\SetCourseField{semesters}{Autumn, Spring}
\SetCourseField{description}{This course is the first in a two semester sequence for teachers of elementary and middle grade students. This course focuses on concepts of numbers and arithmetic operations, including modern and historical perspectives.}
\SetCourseField{prereq}{A grade of C$-$ or above in 1075; or credit for 1074, 75, or 104; or Math Placement Level R or above; or ACT math subscore of 22 or higher that is less than 2 years old.}
\SetCourseField{exclusions}{Not open to students with credit for 106.}

\begin{document}

\CatalogDescription
\PrerequisitesAndExclusions

\begin{Purpose}
This course covers the concepts of whole numbers (positive and negative), place value (base-ten and alternate bases), decimals, and fractions. Some content on irrational numbers appears at the end, and this is extended in Algebra and coordinate geometry for teachers (2137). The four arithmetic operations are covered both conceptually and algorithmically. Attention is given to ensuring that students can perform the algorithms correctly and explain why they give accurate answers. Lastly, the course covers the concepts of proportions and how they are related both to multiplication/division and to fractions. Factors, divisibility, and some elementary number theory complete the course.
\end{Purpose}

\begin{Text}
Mathematics for Elementary Teachers, with Activity Manual, 4th Edition, by Sybilla Beckmann, Pearson, ISBN for the package is 9780321836715 (loose-leaf).
\end{Text}

\begin{TopicsList}
\item Counting numbers, decimals
\item Meaning of fractions
\item Meaning of addition and subtraction
\item Meaning of multiplication
\item Multiplying fractions, decimals, integers
\item Meaning of division
\item Dividing fractions, decimals, integers
\item Meaning of ratios, rates, proportions
\item Greatest common divisor, least common multiple
\item Rational and irrational numbers
\end{TopicsList}

\end{document}
